\documentclass[11pt,a4paper]{article}
\usepackage{curso}
\begin{document}
\portada{Lección}{11.2}{Torx: programación estocástica diferenciable}

\noindent\textbf{Unidad 11.}\quad\textbf{Objetivos:} circuitos estocásticos parametrizados (PSC); la identidad de desplazamiento logarítmico y los estimadores de
gradiente; compuertas derivadas de generadores y el error de componerlas por aristas.
\textbf{Prerrequisitos:} 2.8, 11.1.

\section{Circuitos estocásticos parametrizados}
Verdon et al.\ (agosto de 2026): un PSC es una secuencia de capas de compuertas estocásticas locales sobre cables con tipo (binarios, categóricos, continuos); torx es la
biblioteca en JAX que los construye, ejecuta (simulador de vector de estado, por muestras, o con aleatoriedad de hardware) y deriva. Compuertas de mezcla
$G(\theta)=(1-\sigma(\theta))I+\sigma(\theta)B$: $\mathrm{PNOT}(p)=\left(\begin{smallmatrix}1-p&p\\p&1-p\end{smallmatrix}\right)$, PSWAP, PCNOT. Circuito:
$\rho_{\text{out}}=G_k\cdots G_1\rho_{\text{in}}$.

\textbf{Nota de evidencia.} torx no está instalado en este entorno; los números de esta lección provienen de una reimplementación mínima en NumPy escrita a partir de
las definiciones del artículo (\texttt{datos/psc\_minimo.py}) \evS.

\section{Entrenamiento}
PNOT$\cdot$PSWAP$\cdot$PCNOT sobre $|00\rangle$, objetivo $(\tfrac12,0,0,\tfrac12)$, Adam 3000 pasos: pérdida $0{,}33\to2\cdot10^{-5}$; distribución final
$(0{,}498;\,0{,}002;\,0{,}002;\,0{,}498)$.

\section{Gradientes desde ejecuciones}
$\partial_\theta G=\sigma(1-\sigma)(B-I)=G(\theta)-G(\theta^-)$ con $\sigma(\theta^-)=\sigma(\theta)^2$, es decir $\theta^-=-\ln((1+e^{-\theta})^2-1)$ (verificado con error
$\sim10^{-11}$). Por linealidad, $\partial_{\theta_j}J=\mathbb E_{\theta}[f]-\mathbb E_{\theta_j^-}[f]$: diferencia de dos promedios de ejecuciones. REINFORCE:
$\partial J=\mathbb E[f(z^{(k)})\,\partial\log G_j(z^{(j)}\mid z^{(j-1)})]$, requiere la trayectoria interna.

Con $S=1000$ ejecuciones y 400 repeticiones: exacto $-0{,}0428$; desplazamiento $-0{,}0420\pm0{,}021$; REINFORCE $-0{,}0424\pm0{,}010$. Ambos insesgados; aquí
REINFORCE tuvo menos varianza. La ventaja del desplazamiento es que solo requiere leer la salida.

\section{Compuertas derivadas de generadores}
$\mathrm{PIsing}(\theta)=e^{\Delta t\,Q}$ con tasas de Glauber $Q_{ab}=\sigma(-\beta(E(s_a)-E(s_b)))$. Con $J=1$, $\beta=2$, desde la uniforme: estados alineados
$0{,}294$ ($\Delta t=0{,}1$), $0{,}458$ (1), $0{,}491$ (5) $=$ Boltzmann. Componer PIsing por aristas en un anillo de 8 espines ($\beta=1{,}5$, $J=1$) no da el Ising del
anillo: distancia de variación total estacionaria $0{,}84$, frente a Gibbs cromático exacto. El artículo reporta el mismo efecto ($0{,}402$ en su configuración) y usa
compuertas de Gibbs por sitio.

\section{Hardware}
El artículo ejecuta un Ising de 8 sitios con aleatoriedad de pbits calibrados del chip X0 (plataforma XTR-0, error de calibración $\sim5\,\%$, relajación $\sim100$ ns),
con Metropolis--Hastings y muestreo por importancia, y obtiene distancias consistentes con la referencia en software \evP.

\begin{aplicacion}
Paseos aleatorios, difusión discreta, redes estocásticas en grafos, difusión con saltos; mismo circuito lógico con aleatoriedad de software o hardware.
\end{aplicacion}

\begin{error}
\begin{itemize}[nosep]
\item Componer equilibrios locales esperando el global. \quad REINFORCE sin acceso a estados internos.
\item Probabilidades en 0 o 1 (puntuación divergente). \quad Ignorar la calibración de los pbits reales.
\end{itemize}
\end{error}

\begin{thebibliography}{9}\small
\bibitem{psc} G. Verdon et al., A framework for stochastic differentiable programming, arXiv:2608.01612 (2026).
\bibitem{torx} Extropic, ejemplos de torx (carpeta \texttt{torx/} del repositorio), 2026.
\bibitem{reinforce} R. J. Williams, Simple statistical gradient-following algorithms, \emph{Mach. Learn.} 8, 229 (1992).
\end{thebibliography}
\end{document}
